% 图 3-3 单口共享 SRAM 上加载器写、DMA 读与门铃的时序示意（未按实测比例）
\documentclass[border=4pt]{standalone}
% 所有框图共用的样式（standalone + TikZ，xelatex/tectonic 编译）
\usepackage{fontspec}
\usepackage{xeCJK}
\setCJKmainfont{Noto Sans CJK SC}[BoldFont=Noto Sans CJK SC Bold]
\setCJKsansfont{Noto Sans CJK SC}[BoldFont=Noto Sans CJK SC Bold]
\renewcommand{\familydefault}{\sfdefault}
\usepackage{tikz}
\usetikzlibrary{positioning,arrows.meta,fit,calc,backgrounds,shapes.geometric,decorations.pathreplacing,patterns}
% 统一配色：灰蓝主色，橙色强调，紫色辅助
\definecolor{cBlue}{HTML}{2E6DA4}
\definecolor{cBlueL}{HTML}{DCE7F2}
\definecolor{cBlueM}{HTML}{A9C3DD}
\definecolor{cOrange}{HTML}{C9732E}
\definecolor{cOrangeL}{HTML}{F6E3D2}
\definecolor{cPurple}{HTML}{8A5BB0}
\definecolor{cPurpleL}{HTML}{ECE3F4}
\definecolor{cGray}{HTML}{5B6670}
\definecolor{cGrayL}{HTML}{EEF0F2}
\definecolor{cInk}{HTML}{25303B}
\tikzset{
  >={Stealth[length=2.2mm,width=1.6mm]},
  every node/.style={font=\small, text=cInk},
  blk/.style={draw=cBlue, fill=cBlueL, line width=0.6pt, rounded corners=1.5pt, align=center, minimum height=7mm, inner sep=3pt},
  blkO/.style={blk, draw=cOrange, fill=cOrangeL},
  blkP/.style={blk, draw=cPurple, fill=cPurpleL},
  blkG/.style={blk, draw=cGray, fill=cGrayL},
  blkW/.style={blk, draw=cGray, fill=white},
  grp/.style={draw=cGray, dashed, line width=0.5pt, rounded corners=3pt, inner sep=6pt},
  grpS/.style={draw=cBlue, line width=0.8pt, rounded corners=3pt, inner sep=6pt},
  lbl/.style={font=\footnotesize, text=cGray, inner sep=1pt},
  ttl/.style={font=\small\bfseries, text=cInk, inner sep=1pt},
  ar/.style={->, draw=cInk, line width=0.6pt},
  arB/.style={->, draw=cBlue, line width=0.9pt},
  arO/.style={->, draw=cOrange, line width=0.9pt},
  arD/.style={->, draw=cGray, line width=0.6pt, dashed},
  bus/.style={draw=cGray, line width=2.4pt},
}

\begin{document}
\begin{tikzpicture}[x=1.15cm, y=1cm,
  lw/.style={draw=white, fill=cGray!80, line width=0.3pt, minimum height=0.55cm, anchor=south west, inner sep=0},
  dr/.style={draw=white, fill=cBlue!75, line width=0.3pt, minimum height=0.55cm, anchor=south west, inner sep=0},
  eg/.style={draw=cBlue, fill=cBlueL, line width=0.5pt, minimum height=0.55cm, anchor=south west, inner sep=0, font=\footnotesize},
  ea/.style={draw=cOrange, fill=cOrange!25, line width=0.5pt, minimum height=0.55cm, anchor=south west, inner sep=0, font=\footnotesize},
  db/.style={fill=cPurple, draw=none, isosceles triangle, isosceles triangle apex angle=60, shape border rotate=270, minimum size=2.2mm, inner sep=0},
]
% 行标签
\node[anchor=east, font=\footnotesize] at (-0.15,1.775) {共享 SRAM 端口};
\node[anchor=east, font=\footnotesize] at (-0.15,0.275) {KSE 引擎};
% 端口：首批只有加载器写
\node[lw, minimum width=1.15cm*1.2] at (0,1.5) {};
\node[font=\scriptsize, text=white] at (0.6,1.775) {首批};
% 交替段 1：DMA 读与加载器写 1:1 交替
\foreach \i in {0,...,15} {
  \pgfmathsetmacro{\x}{1.2+0.25*\i}
  \ifodd\i \node[lw, minimum width=1.15cm*0.25] at (\x,1.5) {}; \else \node[dr, minimum width=1.15cm*0.25] at (\x,1.5) {}; \fi
}
% 只有加载器：引擎处于注意力 / LayerNorm 阶段
\node[lw, minimum width=1.15cm*2.0] at (5.2,1.5) {};
\node[font=\scriptsize, text=white] at (6.2,1.775) {加载器独占端口};
% 交替段 2
\foreach \i in {0,...,11} {
  \pgfmathsetmacro{\x}{7.2+0.25*\i}
  \ifodd\i \node[lw, minimum width=1.15cm*0.25] at (\x,1.5) {}; \else \node[dr, minimum width=1.15cm*0.25] at (\x,1.5) {}; \fi
}
% 门铃：每写完 1,024 字发布一次写指针
\foreach \x in {1.2,3.2,5.2,6.2,7.2,9.2} \node[db] at (\x,2.2) {};
\node[font=\scriptsize, text=cPurple, anchor=west] at (9.35,2.2) {WPTR 门铃};
\draw[decorate, decoration={brace, amplitude=4pt}, draw=cGray] (1.2,2.45) -- (3.2,2.45) node[midway, above=4pt, lbl] {交替时 1,024 字约 2,048 拍};
\draw[decorate, decoration={brace, amplitude=4pt}, draw=cGray] (5.2,2.45) -- (6.2,2.45) node[midway, above=4pt, lbl] {约 1,024 拍};
% 引擎
\node[ea, minimum width=1.15cm*1.2] at (0,0) {等待首批};
\node[eg, minimum width=1.15cm*4.0] at (1.2,0) {矩阵乘：消费参数，约 0.5 字/拍};
\node[ea, minimum width=1.15cm*2.0] at (5.2,0) {注意力 / LN};
\node[eg, minimum width=1.15cm*3.0] at (7.2,0) {矩阵乘：消费参数};
\draw[arD] (1.2,1.5) -- (1.2,0.55);
\node[lbl, anchor=north] at (1.2,-0.05) {\scriptsize 放行：A + 4 $\le$ WPTR};
% 时间轴
\draw[->, draw=cInk, line width=0.6pt] (0,-0.75) -- (10.5,-0.75) node[right, font=\footnotesize] {周期};
% 图例
\node[lw, minimum width=4mm, minimum height=3mm, anchor=center] (l0) at (0.2,3.3) {};
\node[anchor=west, font=\scriptsize] at (l0.east) {加载器写（含门铃）};
\node[dr, minimum width=4mm, minimum height=3mm, anchor=center] (l1) at (3.0,3.3) {};
\node[anchor=west, font=\scriptsize] at (l1.east) {DMA 读（优先）};
\node[ea, minimum width=4mm, minimum height=3mm, anchor=center] (l2) at (5.5,3.3) {};
\node[anchor=west, font=\scriptsize] at (l2.east) {引擎不读参数的阶段};
\node[anchor=west, font=\scriptsize, text=cGray] at (0.05,-1.2) {示意图，未按实测比例；端口每拍至多完成一笔事务};
\end{tikzpicture}
\end{document}
